\chapter*{Apresentação}
\addcontentsline{toc}{chapter}{Apresentação}
\markboth{Apresentação}{Apresentação}
Esta apostila reúne os fundamentos e as tendências dos \textbf{Sistemas Supervisórios e SCADA} (\emph{Supervisory Control and Data Acquisition}) aplicados à automação industrial. O material foi desenvolvido para oferecer uma progressão didática --- do conceito à implementação --- com linguagem acessível, diagramas explicativos e estudos de caso reais.

Os 20 capítulos estão organizados em quatro blocos:

\begin{itemize}
\tightlist
\item
  \textbf{Fundamentos (capítulos 1 a 4):} automação industrial, pirâmide de automação e níveis hierárquicos, sistemas supervisórios (SCADA e HMI), arquitetura e componentes.
\item
  \textbf{Comunicação, integração e segurança (capítulos 5 a 9 e 13):} redes e protocolos industriais, IIoT, Node-RED, cibersegurança, Cloud SCADA e Digital Twin, OPC UA, MQTT e Sparkplug B.
\item
  \textbf{Dados, alarmes e visualização (capítulos 10 a 12 e 14 a 16):} tags, variáveis, estados e eventos; aquisição e tratamento de dados; historiadores e séries temporais; gerenciamento de alarmes; dashboards; ThingsBoard.
\item
  \textbf{Borda, nuvem e projeto (capítulos 17 a 20):} edge computing, integração SCADA + CLP + IoT + Cloud, projeto e implantação de um sistema supervisório, IA e analítica industrial.
\end{itemize}

As normas aparecem onde mudam a decisão de engenharia, não como apêndice teórico: o \textbf{gerenciamento de alarmes} segue ANSI/ISA-18.2, IEC 62682, EEMUA 191 e NAMUR NA 102, e a \textbf{cibersegurança industrial} se apoia na série IEC 62443.

Cada capítulo abre com objetivos de aprendizagem e fecha com questões de revisão --- 159 no total. O corpo do texto traz 47 diagramas Mermaid, 64 tabelas legendadas e 79 posições de figura, 31 delas com reproduções do material de origem (telas de supervisão, arquiteturas de rede, painéis de alarmes, tendências e dashboards). Todos os 20 capítulos têm \textbf{estudo de caso}, com cenário, diagnóstico, decisão de engenharia e limitações do arranjo; os capítulos 10 a 13 trazem \textbf{atividades práticas}; e o \textbf{glossário técnico} em apêndice resolve a terminologia durante a leitura.

Ao final do percurso, o leitor deve ser capaz de especificar, integrar e operar um sistema supervisório completo --- e de justificar cada escolha de arquitetura, protocolo, base de dados e política de alarmes.

\begin{center}\rule{0.5\linewidth}{0.5pt}\end{center}
